\begingroup\itshape
\section{\textit{Introducción}}

En octubre de 2026, Colombia formula el Plan Nacional de Desarrollo 2026–2030. El DNP abrió la recepción de propuestas de la academia, las empresas y los territorios. El Centro Javeriano de Competitividad presenta diez propuestas para elevar la productividad de las empresas, aumentar los ingresos de los trabajadores y mejorar los servicios públicos. \href{https://dnp.gov.co/Prensa\_/Noticias/Paginas/aqui-la-patria-se-construye-unida-el-dnp-lanza-el-buzon-digital-para-que-todos-los-colombianos-escriban-su-propuesta.aspx}{DNP abrió la recepción de propuestas}.

La inteligencia artificial (IA) utiliza datos para identificar patrones, hacer predicciones o generar contenido. La generación de textos, imágenes u otros contenidos se conoce como IA generativa. Estas herramientas pueden apoyar tareas como revisar documentos, detectar defectos o prever inventarios. Su conveniencia depende de los datos disponibles, los errores y el costo de implementación y revisión.

Las propuestas atienden problemas de acceso a asistencia técnica, financiación de investigación, formación y coordinación entre entidades. El objetivo es que empresas y Estado escojan herramientas que mejoren sus procesos y comprueben sus resultados antes de ampliar su uso. La propuesta sobre ahorro pensional tiene un alcance mayor que la IA y busca mejorar la financiación empresarial mediante capital.

La primera amplía la asistencia tecnológica a mipymes para diagnosticar procesos, ordenar datos, implementar herramientas y presupuestar cuánto costará mantenerlas después del apoyo.

La segunda promueve alianzas estables de empresas y universidades para investigar problemas productivos. Los acuerdos separarán el sostenimiento de equipos científicos y laboratorios de la financiación de cada proyecto. Definirán cuánto dinero pagará cada empresa, qué personal, materiales o equipos aportará y quién podrá utilizar los resultados.

La tercera combina formación docente y evaluación de usos educativos de IA con formación laboral por ocupación. Los bonos permitirán pagar cursos con evaluaciones prácticas de las tareas que el participante debe aprender. Se hará seguimiento del empleo y los ingresos de quienes reciban formación laboral.

La cuarta cofinancia proyectos conjuntos de desarrolladores de IA y empresas usuarias. Cada etapa tendrá entregas y pruebas acordadas, como datos listos para trabajar, una herramienta que funcione o una comparación con el procedimiento actual. Su cumplimiento determinará la continuidad del apoyo.

La quinta adapta la asistencia a firmas contables, jurídicas y administrativas. Las herramientas deberán mantener un registro de documentos y fuentes utilizadas, permitir la revisión del profesional responsable y proteger la información de los clientes.

La sexta utiliza los mecanismos de pruebas supervisadas de la SFC para evaluar aplicaciones financieras y aseguradoras. Cada prueba limitará duración y usuarios expuestos, medirá errores y efectos sobre los clientes y permitirá detener la aplicación.

La séptima ofrece acceso temporal a capacidad de procesamiento informático cuando su costo limite un proyecto viable. Se compararán servicios contratados y equipos compartidos según la demanda y el costo total.

La octava elimina la garantía de rentabilidad mínima del ahorro pensional y fija en \$2.000 millones el mínimo inicial de capital para administrar pensiones. Busca mejorar la asignación del riesgo y facilitar financiación mediante capital, con gestión prudente y protección de los afiliados.

La novena ofrece herramientas de búsqueda de documentos, clasificación de solicitudes y preparación de borradores que puedan utilizar varias entidades públicas. Se medirán tiempos, errores y costos, incluida la revisión del funcionario. Cada entidad conservará la responsabilidad sobre sus decisiones y datos.

La décima reúne los proyectos de IA que recibirán prioridad y financiación durante el cuatrienio. Cada proyecto deberá identificar qué cambiará en tiempos, errores, costos o calidad. El DNP coordinará el seguimiento y las entidades sectoriales responderán por la ejecución.
\par\endgroup
